\documentclass[10pt,a4paper]{amsart}
\usepackage[margin=19mm]{geometry}
\usepackage{amsmath,amssymb,mathtools}
\usepackage[scr=rsfs,cal=euler]{mathalfa}
\usepackage{xfrac}
\usepackage[unicode,hidelinks,bookmarks=false]{hyperref}
\newcommand{\ie}{i.e.\ }
\newcommand{\cf}{cf.\ }
\newcommand{\resp}{respectively\ }
\newcommand{\Subsection}{\subsection}
\providecommand{\nobreakdash}{\nobreak\hbox{-}}
\setlength{\emergencystretch}{2em}
\multlinegap0pt
\usepackage{amssymb}
\usepackage{float}
\newtheorem{lem}{Lemma}
\newcommand{\F}{{\mathcal{F}}}
\newcommand{\beas}{\begin{eqnarray*}}
\newcommand{\eeas}{\end{eqnarray*}}
\renewcommand{\epsilon}{\varepsilon}
\newcommand{\ol}{\overline}
\title{Generalized Bohr inequalities for K-quasiconformal harmonic mappings and their applications}
\author{Raju Biswas, Rajib Mandal}
\date{Source: arXiv:2411.01837v1 (CC BY 4.0)}
\begin{document}
\maketitle

\noindent\emph{Source and licence.} Generalized Bohr inequalities for K-quasiconformal harmonic mappings and their applications. Version arXiv:2411.01837v1 (CC BY 4.0), distributed by arXiv under the Creative Commons Attribution 4.0 International licence, http://creativecommons.org/licenses/by/4.0/, as recorded in the arXiv licence metadata for this exact version and shown on the abstract page https://arxiv.org/abs/2411.01837v1. Full licence notice: \texttt{licenses/arxiv-2411.01837-CC-BY-4.0.txt}.\par\smallskip

\noindent\textit{Editorial scope.} Counted unit: the article's lem lem7 (source0.tex lines 245--248). The article's complete original proof of it is delivered with it (source0.tex lines 249--263, one \texttt{proof} environment of 15 lines). No mathematics is written, repaired or re-derived: the statement and the proof are the article's own lines, in the article's own notation, and no retained line is substituted or re-flowed.  Retained uncounted as the source-local context the proof needs: lines 227--237.  Nothing was omitted from any retained span, so the package carries no omission record.  The retained lines cite 2 works, whose entries are transcribed verbatim from the article's own bibliography source in the e-print and delivered with short editorial labels recorded in the item's reference mapping.  No retained line contains a figure environment, an image-inclusion command or drawing code, so no image is delivered and the package declares no asset dependency.  Editorial formatting only, in addition: an AMS \texttt{amsart} preamble that restates the article's own declarations and macros used by the retained lines, namely the environments \texttt{lem} on separate counters, together with the standard packages those lines need.  The retained environments print in this document as Lemma 1 in order of appearance, whereas the article prints the same items with the numbering of its own full document.  The complete native source of the article as distributed in the arXiv e-print for this exact version is bundled unchanged as \texttt{sources/167915/source0.tex}.
\begin{lem}\label{lem3}\cite{C1940}\cite[Lemma B]{PVW2020} Suppose that $f\in\mathcal{B}(\mathbb{D},\mathbb{C})$ and $f(z) = \sum_{k=0}^\infty a_kz^k$. Then the following inequalities hold.
\begin{enumerate}
\item[(i)] $\vert a_{2k+1}\vert\leq 1-\vert a_0\vert^2-\cdots-\vert a_k\vert^2$, $k=0,1,\ldots$
\item[(ii)] $\vert a_{2k}\vert\leq 1-\vert a_0\vert^2-\cdots-\vert a_{k-1}\vert^2-\frac{\vert a_k\vert^2}{1+\vert a_0\vert}$, $k=1,2,\ldots$.
\end{enumerate}
Further, to have equality in $(i)$ it is necessary that $f$ is a rational function of the form 
\beas f(z)=\frac{a_0+a_1z+\cdots+a_nz^n+\epsilon z^{2n+1}}{1+(\ol{a_n}z^n+\cdots+\ol{a_0}z^{2n+1})\epsilon},\quad|\epsilon|=1,\eeas
and to have the equality in $(ii)$ it is necessary that $f$ is a rational function of the form 
\beas f(z)=\frac{a_0+a_1z+\cdots+(a_n/(1+|a_0|))z^n+\epsilon z^{2n+1}}{1+((\ol{a_n}/(1+|a_0|))z^n+\cdots+\ol{a_0}z^{2n})\epsilon},\quad|\epsilon|=1,\eeas
where the term $a_0\ol{a_n}^2\epsilon$ is non-negative real.
\end{lem}

\begin{lem}\label{lem7} Suppose $f(z)=\sum_{n=0}^\infty a_nz^n$ is analytic in $\mathbb{D}$ with $|f(z)|\leq1$ and $\{\psi_n(r)\}_{n=1}^\infty \in\F$. Then, we have the following inequality:
\beas \sum_{n=2}^\infty |a_n|\psi_n(r)+\sum_{n=1}^\infty |a_n|^2\left(\frac{\psi_{2n}(r)}{1+|a_0|}+\Psi_{2n+1}(r)\right)\leq \left(1-|a_0|^2\right)\Psi_2(r),\eeas
where $\Psi_t(r)=\sum_{k=t}^\infty \psi_k(r)$.
\end{lem}

\begin{proof} In view of \textrm{Lemma \ref{lem3}}, we have
\beas\sum_{n=2}^\infty |a_n|\psi_n(r)&=&\sum_{n=1}^\infty |a_{2n}|\psi_{2n}(r)+\sum_{n=1}^\infty |a_{2n+1}|\psi_{2n+1}(r)\\[2mm]
&\leq &\sum_{n=1}^\infty \left(1-|a_0|^2-\sum_{k=1}^{n-1}|a_k|^2-\frac{|a_n|^2}{1+|a_0|}\right)\psi_{2n}(r)\\[2mm]
&&+\sum_{n=1}^\infty \left(1-|a_0|^2-\sum_{k=1}^{n}|a_k|^2\right)\psi_{2n+1}(r)\\[2mm]
&=&\left(1-|a_0|^2\right)\sum_{n=2}^\infty \psi_n(r)-\sum_{n=1}^\infty\frac{|a_n|^2}{1+|a_0|}\psi_{2n}(r)\\[2mm]
&&-\left(\sum_{n=1}^\infty\left(\sum_{k=1}^{n-1}|a_k|^2\right) \psi_{2n}(r)+\sum_{n=1}^\infty\left(\sum_{k=1}^{n}|a_k|^2\right) \psi_{2n+1}(r)\right).\eeas
It is evident that 
\beas &&\sum_{n=1}^\infty\left(\sum_{k=1}^{n-1}|a_k|^2\right) \psi_{2n}(r)+\sum_{n=1}^\infty\left(\sum_{k=1}^{n}|a_k|^2\right) \psi_{2n+1}(r)\\[2mm]
&=&|a_1|^2\sum_{k=3}^\infty \psi_k(r)+|a_2|^2\sum_{k=5}^\infty \psi_k(r)+|a_3|^2\sum_{k=7}^\infty \psi_k(r)+\cdots\\[2mm]
&=&\sum_{n=1}^\infty |a_n|^2\left(\sum_{k=2n+1}^\infty \psi_k(r)\right).\eeas
Therefore, we have 
\beas \sum_{n=2}^\infty |a_n|\psi_n(r)&\leq&\left(1-|a_0|^2\right)\sum_{n=2}^\infty \psi_n(r)-\sum_{n=1}^\infty\frac{|a_n|^2}{1+|a_0|}\psi_{2n}(r)-\sum_{n=1}^\infty |a_n|^2\Psi_{2n+1}\\[2mm]
&=&\left(1-|a_0|^2\right)\sum_{n=2}^\infty \psi_n(r)-\sum_{n=1}^\infty|a_n|^2\left(\frac{\psi_{2n}(r)}{1+|a_0|}+\Psi_{2n+1}(r)\right).\eeas
This completes the proof.
\end{proof}

\begin{thebibliography}{99}
\bibitem[C1940]{C1940} {\sc F. Carlson}, Sur les coefficients d'une fonction born\'ee dans le cercle unit\'e (French), {\it Ark. Mat. Astr. Fys.} {\bf 27A}(1) (1940), 8 pp.
\bibitem[PVW202]{PVW2020} {\sc S. Ponnusamy, R. Vijayakumar} and {\sc K. -J. Wirths}, New inequalities for the coefficients of unimodular bounded functions, {\it Results Math} {\bf 75} (2020), 107.
\end{thebibliography}
\end{document}
